\documentclass{beamer}
\usetheme{Madrid}
\usepackage{amsmath}

\title{Measuring Production and Income}
\subtitle{ECON 101 Macroeconomics}
\author{Instructor: Muhebullah Karimzada}

\date{Winter 2025}

\begin{document}

\begin{frame}
\titlepage
\end{frame}

% Problem 1
\begin{frame}{Problem 1: Calculating GDP}
A country produces the following goods and services in a year:
\begin{itemize}
    \item 100 units of cars at \$20,000 each
    \item 500 units of bread at \$5 each
    \item 200 units of smartphones at \$500 each
    \item Services worth \$1,000,000
\end{itemize}
Calculate the Gross Domestic Product (GDP) of the country.
\end{frame}

\begin{frame}{Solution to Problem 1}
GDP is the total market value of all final goods and services produced within a country in a given period. It is calculated as:
\[
\text{GDP} = \sum (\text{Quantity} \times \text{Price})
\]
\[
\text{GDP} = (100 \times 20,000) + (500 \times 5) + (200 \times 500) + 1,000,000
\]
\[
\text{GDP} = 2,000,000 + 2,500 + 100,000 + 1,000,000 = 3,102,500
\]
\textbf{Answer:} GDP = \textbf{\$3,102,500}
\end{frame}

% Problem 2
\begin{frame}{Problem 2: GDP vs. GNP}
A country has the following data:
\begin{itemize}
    \item GDP = \$10 billion
    \item Income earned by domestic residents from abroad = \$2 billion
    \item Income earned by foreigners in the domestic economy = \$1 billion
\end{itemize}
Calculate the Gross National Product (GNP).
\end{frame}

\begin{frame}{Solution to Problem 2}
GNP is the total income earned by a country's residents, including income from abroad. It is calculated as:
\[
\text{GNP} = \text{GDP} + \text{Income from abroad} - \text{Income to foreigners}
\]
\[
\text{GNP} = 10 + 2 - 1 = 11
\]
\textbf{Answer:} GNP = \textbf{\$11 billion}
\end{frame}

% Problem 3
\begin{frame}{Problem 3: Net Domestic Product (NDP)}
A country has the following data:
\begin{itemize}
    \item GDP = \$15 billion
    \item Depreciation = \$2 billion
\end{itemize}
Calculate the Net Domestic Product (NDP).
\end{frame}

\begin{frame}{Solution to Problem 3}
NDP is GDP minus depreciation (the wear and tear on capital goods). It is calculated as:
\[
\text{NDP} = \text{GDP} - \text{Depreciation}
\]
\[
\text{NDP} = 15 - 2 = 13
\]
\textbf{Answer:} NDP = \textbf{\$13 billion}
\end{frame}

% Problem 4
\begin{frame}{Problem 4: National Income}
A country has the following data:
\begin{itemize}
    \item NDP = \$12 billion
    \item Indirect taxes = \$1 billion
    \item Subsidies = \$0.5 billion
\end{itemize}
Calculate the National Income.
\end{frame}

\begin{frame}{Solution to Problem 4}
National Income is NDP minus indirect taxes plus subsidies. It is calculated as:
\[
\text{National Income} = \text{NDP} - \text{Indirect taxes} + \text{Subsidies}
\]
\[
\text{National Income} = 12 - 1 + 0.5 = 11.5
\]
\textbf{Answer:} National Income = \textbf{\$11.5 billion}
\end{frame}

% Problem 5
\begin{frame}{Problem 5: Personal Income}
A country has the following data:
\begin{itemize}
    \item National Income = \$20 billion
    \item Corporate taxes = \$3 billion
    \item Undistributed corporate profits = \$2 billion
    \item Social security contributions = \$1 billion
    \item Transfer payments = \$4 billion
\end{itemize}
Calculate the Personal Income.
\end{frame}

\begin{frame}{Solution to Problem 5}
Personal Income is National Income minus corporate taxes, undistributed profits, and social security contributions, plus transfer payments. It is calculated as:
\[
\text{Personal Income} = \text{National Income} - \text{Corporate taxes} - \text{Undistributed profits} - \text{Social security} + \text{Transfer payments}
\]
\[
\text{Personal Income} = 20 - 3 - 2 - 1 + 4 = 18
\]
\textbf{Answer:} Personal Income = \textbf{\$18 billion}
\end{frame}

% Problem 6
\begin{frame}{Problem 6: Disposable Income}
A country has the following data:
\begin{itemize}
    \item Personal Income = \$25 billion
    \item Personal taxes = \$5 billion
\end{itemize}
Calculate the Disposable Income.
\end{frame}

\begin{frame}{Solution to Problem 6}
Disposable Income is Personal Income minus personal taxes. It is calculated as:
\[
\text{Disposable Income} = \text{Personal Income} - \text{Personal taxes}
\]
\[
\text{Disposable Income} = 25 - 5 = 20
\]
\textbf{Answer:} Disposable Income = \textbf{\$20 billion}
\end{frame}

% Problem 7
\begin{frame}{Problem 7: Real GDP}
A country has the following data:
\begin{itemize}
    \item Nominal GDP in Year 1 = \$1,000 billion
    \item Nominal GDP in Year 2 = \$1,200 billion
    \item Price index in Year 1 = 100
    \item Price index in Year 2 = 110
\end{itemize}
Calculate the Real GDP in Year 2 using Year 1 as the base year.
\end{frame}

\begin{frame}{Solution to Problem 7}
Real GDP adjusts Nominal GDP for inflation using the price index. It is calculated as:
\[
\text{Real GDP} = \frac{\text{Nominal GDP}}{\text{Price index}} \times 100
\]
\[
\text{Real GDP} = \frac{1,200}{110} \times 100 = 1,090.91
\]
\textbf{Answer:} Real GDP = \textbf{\$1,090.91 billion}
\end{frame}

% Problem 8
\begin{frame}{Problem 8: GDP Deflator}
A country has the following data:
\begin{itemize}
    \item Nominal GDP = \$1,500 billion
    \item Real GDP = \$1,200 billion
\end{itemize}
Calculate the GDP deflator.
\end{frame}

\begin{frame}{Solution to Problem 8}
The GDP deflator measures the level of prices in the economy. It is calculated as:
\[
\text{GDP Deflator} = \frac{\text{Nominal GDP}}{\text{Real GDP}} \times 100
\]
\[
\text{GDP Deflator} = \frac{1,500}{1,200} \times 100 = 125
\]
\textbf{Answer:} GDP Deflator = \textbf{125}
\end{frame}

% Problem 9
\begin{frame}{Problem 9: Per Capita GDP}
A country has the following data:
\begin{itemize}
    \item GDP = \$500 billion
    \item Population = 50 million
\end{itemize}
Calculate the GDP per capita.
\end{frame}

\begin{frame}{Solution to Problem 9}
GDP per capita is GDP divided by the population. It is calculated as:
\[
\text{GDP per capita} = \frac{\text{GDP}}{\text{Population}}
\]
\[
\text{GDP per capita} = \frac{500,000,000,000}{50,000,000} = 10,000
\]
\textbf{Answer:} GDP per capita = \textbf{\$10,000}
\end{frame}

% Problem 10
\begin{frame}{Problem 10: Expenditure Approach to GDP}
A country has the following data:
\begin{itemize}
    \item Consumption = \$400 billion
    \item Investment = \$150 billion
    \item Government spending = \$200 billion
    \item Exports = \$100 billion
    \item Imports = \$50 billion
\end{itemize}
Calculate the GDP using the expenditure approach.
\end{frame}

\begin{frame}{Solution to Problem 10}
The expenditure approach calculates GDP as the sum of consumption, investment, government spending, and net exports (exports minus imports). It is calculated as:
\[
\text{GDP} = C + I + G + (X - M)
\]
\[
\text{GDP} = 400 + 150 + 200 + (100 - 50) = 800
\]
\textbf{Answer:} GDP = \textbf{\$800 billion}
\end{frame}

\end{document}